% TOP_PROBLEM: {"id": 30006865, "title": "Arnold Chord Conjecture in Higher Dimensions", "category_id": 6, "category": {"id": 6, "name": "geometry", "display_name": "Geometry", "description": "Euclidean and non-Euclidean geometry, geometric structures.", "slug": "geometry", "order_index": 6, "created_at": "2026-07-31T15:26:25.671Z"}, "status": "open"}
% FIRST_POSTED: 2026-09-27
% !TeX program = lualatex
% ATTEMPT_STATUS: unresolved
% Model: GPT-6 Astra Ultra; continuation dated 27 September 2026.
\documentclass[11pt]{article}
\usepackage[margin=25mm]{geometry}
\usepackage{fontspec}
\setmainfont{Latin Modern Roman}
\usepackage{amsmath,amssymb,mathtools}
\usepackage{unicode-math}
\setmathfont{Latin Modern Math}
\usepackage{xurl,hyperref}
\hypersetup{colorlinks=true,urlcolor=blue}
\setcounter{secnumdepth}{0}
\setlength{\parindent}{0pt}
\setlength{\parskip}{5pt}
\setlength{\emergencystretch}{3em}
\begin{document}
\section*{\texorpdfstring{Arnold Chord Conjecture in Higher Dimensions}{Arnold Chord Conjecture in Higher Dimensions}}\label{30006865}
\subsection{Definitions and mathematical statement}
A co-oriented contact structure on a closed $(2n+1)$-manifold $Y$ is $\xi=\ker\alpha$ for a one-form with $\alpha\wedge({\,\mathrm{d}}\alpha)^n$ nowhere zero. Its Reeb vector field is uniquely determined by
\[
 \alpha(R_\alpha)=1,\qquad \iota_{R_\alpha}{\,\mathrm{d}}\alpha=0.
\]
A Legendrian submanifold $\Lambda\subset Y$ is a closed $n$-dimensional submanifold with $T\Lambda\subset\xi$. The higher-dimensional chord assertion says that for every such datum with $n\ge2$, there are $T>0$ and $x\in\Lambda$ such that $\varphi_{R_\alpha}^T(x)\in\Lambda$. Thus $t\mapsto\varphi_{R_\alpha}^t(x)$, $0\le t\le T$, is a nonconstant Reeb chord. Its two endpoints need not coincide, and a closed Reeb orbit disjoint from $\Lambda$ would not answer the question.
\par\smallskip\textit{Formulation and concept references:} \href{https://arxiv.org/abs/1004.4319}{[S1]}.

\subsection{Short English statement}
Must every closed Legendrian submanifold in a closed higher-dimensional contact manifold have a nonconstant Reeb trajectory beginning and ending on it?
\subsection{Research attempt}
\textbf{Scope and current-source check.}
The three-dimensional theorem in [S1] does not cover arbitrary closed contact
manifolds of dimension at least five.  Two recent records require particular
care.  The general proof claim [E1] was withdrawn in version 2, dated
17 September 2025; the author's stated issue is inability to exclude a
constant limiting orbit.  It is not used as a theorem here.  The August 2026
preprint [E2] proves a broad class of chord results under contact
non-orderability and symplectization rigidity hypotheses, including specified
prequantization cases.  Its abstract does not assert the universal target
without those hypotheses.  The following arguments give direct subcases and
explain an obstruction to an elementary symplectization proof.

\textbf{The positive-time requirement is geometrically substantive.}
Cartan's formula gives
\[
 \mathcal L_{R_\alpha}\alpha
 =d(\alpha(R_\alpha))+\iota_{R_\alpha}d\alpha=0.
\]
Thus the Reeb flow preserves $\alpha$, and therefore its contact distribution
and the volume form $\alpha\wedge(d\alpha)^n$.  Also $R_\alpha$ never vanishes,
since $\alpha(R_\alpha)=1$.  For a compact embedded Legendrian $\Lambda$, there
is $\tau>0$ such that no Reeb chord from $\Lambda$ to itself has length in
$(0,\tau)$.  To prove this, the map
\[
 F:\Lambda\times\mathbb R\longrightarrow Y,\qquad F(x,t)=\varphi^t(x)
\]
has injective differential at every $(x,0)$: its tangent image is the direct
sum of $T_x\Lambda$ and $\mathbb R R_\alpha(x)$, because $\alpha$ vanishes on
the former and equals one on the latter.  The flow-box theorem gives a local
product neighborhood in which two sufficiently nearby points of $\Lambda$
cannot be joined by a sufficiently short nonzero Reeb segment.  If a uniform
$\tau$ failed, compactness would yield $x_j\to x$, $t_j\downarrow0$ with
$\varphi^{t_j}(x_j)\in\Lambda$, contradicting that local product neighborhood.
This proves a positive lower bound for chord times for this fixed datum, not
existence of a chord or a uniform bound over all contact forms.

\textbf{Periodic flows and an explicit irrational-flow example.}
Whenever $\varphi^{T_0}=\mathrm{id}$ for some $T_0>0$, every nonempty
Legendrian has a chord of length $T_0$.  Its endpoints may be identical.
On the unit sphere in $\mathbb C^{n+1}$, the restriction of
\[
 \alpha_0=\frac12\sum_j(x_j\,dy_j-y_j\,dx_j)
\]
has Reeb field $R=2\sum_j(-y_j\partial_{x_j}+x_j\partial_{y_j})$.
Indeed $\alpha_0(R)=\sum_j|z_j|^2=1$, and contraction with
$d\alpha_0=\sum_j dx_j\wedge dy_j$ vanishes on the sphere's tangent space.
The flow is $z\mapsto e^{2it}z$, with common period $\pi$.  This proves the
chord assertion for every closed Legendrian and this particular form in all
dimensions.  Multiplying the form by a positive constant multiplies the
period by that constant; multiplying it by an arbitrary function does not
preserve this elementary argument.

There is also a useful example without a common period.  On the ellipsoid
\[
 Y_a=\left\{z:\sum_j\frac{|z_j|^2}{a_j}=1\right\},\qquad a_j>0,
\]
the same form has Reeb flow $z_j(t)=e^{2it/a_j}z_j(0)$.
The real locus $\Lambda_a=Y_a\cap\mathbb R^{n+1}$ is a closed Legendrian:
its tangent vectors have $dy_j=0$, while $y_j=0$.  At a coordinate-axis point
$z=\sqrt{a_j}e_j$, the time $T=\pi a_j/2$ reaches
$-\sqrt{a_j}e_j\in\Lambda_a$.  These are distinct endpoints.
The calculation holds even when all frequency ratios are irrational.  It
does not prove the assertion for every Legendrian on an irrational ellipsoid;
the real locus is an additional special feature of this construction.

\textbf{Cotangent fibres: a topological subcase.}
Let $Q$ be a closed Riemannian manifold with $\pi_1(Q)\ne0$.  On its unit
cotangent bundle, the canonical contact form has geodesic Reeb flow.  For a
fixed $q\in Q$, the fibre $\Lambda=S_q^*Q$ is Legendrian.  Select a nontrivial
based homotopy class of loops at $q$.  It contains a shortest rectifiable
loop.  One way to see this is to take a minimizing sequence parametrized with
constant speed on $[0,1]$.  Bounded lengths give equicontinuity; compactness
and Arzel\`a--Ascoli give a uniform limit.  Sufficiently uniformly close
based loops are based-homotopic using short geodesic segments, so the limit
retains the class.  Length is lower semicontinuous, hence minimal.  The
minimum is positive, since a sufficiently short based loop lies in a normal
ball and is contractible.

Replacing any short interior subarc by a minimizing geodesic cannot lower
length, so the minimizing loop is a geodesic on $(0,1)$.  Its endpoint
velocities need not agree.  Lifting this nonconstant geodesic segment to the
unit cotangent bundle gives a Reeb chord from $S_q^*Q$ to itself.  This is
precisely why requiring equal covector endpoints would be unnecessarily
strong.  The proof treats a fibre Legendrian and a Riemannian cosphere form;
it does not extend without argument to an arbitrary contact form or arbitrary
Legendrian in the same contact manifold.

\textbf{What a chord-free flowout would imply.}
Assume, for the purpose of a possible contradiction, that $\Lambda$ has no
positive Reeb chord.  Then $F:\Lambda\times\mathbb R\to Y$ above is injective.
Indeed $\varphi^t(x)=\varphi^s(y)$ with $t>s$ implies
$\varphi^{t-s}(x)=y\in\Lambda$, and $t=s$ implies $x=y$.
Its differential is injective everywhere by flow invariance of the contact
form.  This is an injective immersion, but it need not be a proper embedding
into the compact manifold $Y$.  Compactness therefore gives accumulation,
not an exact intersection at a positive time.

Poincar\'e recurrence is no substitute.  It concerns almost every point for
contact volume, while the $n$-dimensional Legendrian has zero measure in the
$(2n+1)$-dimensional contact manifold.  Even recurrence of a selected point
would only put an orbit arbitrarily close to its initial point, not precisely
back on $\Lambda$.  Thickening $\Lambda$ makes recurrence applicable to a
positive-volume set but loses the exact endpoint condition upon shrinking
unless one has a uniform positive upper bound for the return times.

\textbf{A Lagrangian construction and its exactness obstruction.}
Consider the symplectization $(\mathbb R_s\times Y,d(e^s\alpha))$.
For a simple smooth closed curve $c(\theta)=(s(\theta),t(\theta))$ in the
$(s,t)$-plane, define
\[
 \Psi(x,\theta)=(s(\theta),\varphi^{t(\theta)}(x)).
\]
Under the chord-free assumption this is an embedding of
$\Lambda\times S^1$: equality of images first forces equality of the $s$
coordinates and then, unless the $t$ coordinates agree, a chord; simplicity
of $c$ handles the remaining case.  Compactness of the domain upgrades an
injective immersion to an embedding.  Flow preservation of $\alpha$ and
$\alpha|_{T\Lambda}=0$ imply
\[
 \Psi^*(e^s\alpha)=e^{s(\theta)}t'(\theta)\,d\theta,
 \qquad \Psi^*d(e^s\alpha)=0.
\]
Its dimension is $n+1$, half the symplectization dimension, so it is
Lagrangian.  But the integral of the pulled-back primitive on a circle fibre is
\[
 \oint_c e^s\,dt=\iint_{D_c}e^s\,ds\,dt\ne0
\]
for the positive orientation, by Green's theorem.  It is consequently
\emph{not exact}.  An obstruction to closed exact Lagrangians cannot be applied
to this object.  Conversely the open flowout at fixed $s=0$ has primitive
$dt$, hence is exact, but it is noncompact; truncating it introduces boundary.
Neither alternative yields the needed contradiction by elementary exactness.

This calculation does indicate why symplectic area or displacement estimates
are relevant: the enclosed weighted area is an actual invariant of the
construction.  To turn it into a chord theorem one needs an additional
rigidity mechanism bounding that area, together with all compactness and
nonconstant-limit conclusions.  Such a mechanism is extra input, not a
consequence of the contact equations alone.

\textbf{Verification and outcome.}
The saved diagnostic checks the sphere and ellipsoid Reeb identities at
sample points and evaluates the weighted area identity for rectangles
numerically.  Those checks confirm the normalizations $\pi$ and
$\pi a_j/2$; they do not test arbitrary contact manifolds.  The proofs above
establish periodic-flow and cotangent-fibre subcases and expose the exactness
gap in a natural flowout attack.  No positive upper bound for a general
Legendrian's first chord time is obtained, and the universal higher-dimensional
statement is not proved here.
\subsection{Sources}
\begin{itemize}
\item[S1]Proof of the Arnold chord conjecture in three dimensions I (Hutchings-Taubes, arXiv abstract). Accessed 2026-09-01.
\url{https://arxiv.org/abs/1004.4319}.
\textit{Formulation source.}
\item[E1]Yong-Geun Oh, \emph{Contact instantons and Proofs of Weinstein's conjecture and Arnold's chord conjecture}, arXiv:2412.20731v2, withdrawn 17 September 2025.
\url{https://arxiv.org/abs/2412.20731v2}.
\textit{Withdrawal and stated constant-orbit issue checked; not used as an established theorem.}
\item[E2]Egor Shelukhin, \emph{On orderability and the chord conjecture}, arXiv:2608.10493v1, 11 August 2026.
\url{https://arxiv.org/abs/2608.10493v1}.
\textit{Current primary-source abstract audit; the described result has additional contact and symplectization hypotheses.}
\end{itemize}
\end{document}
